\documentclass[11pt,a4paper]{article}

\usepackage[margin=1in]{geometry}
\usepackage{amsmath,amssymb,amsfonts}
\usepackage{booktabs}
\usepackage{array}
\usepackage{graphicx}
\usepackage{hyperref}
\usepackage{caption}
\usepackage{subcaption}
\usepackage{float}
\usepackage{xcolor}

\hypersetup{
    colorlinks=true,
    linkcolor=blue,
    citecolor=blue,
    urlcolor=blue
}

\graphicspath{{./}{results/}}

\title{\textbf{Jump-Diffusion Schrodinger Bridges for Financial Time-Series Generation}}
\author{
Lizhan Hong \and Sheng Wan \and Haoyu Luo \and Cyprien Kutek \and Dimitri Krawczyk\\
\small Projet Scientifique Collectif, Ecole Polytechnique\\
\small Supervised by Prof. Huyen Pham
}
\date{Academic year 2025--2026}

\begin{document}

\maketitle

\begin{abstract}
This report summarizes the implementation and experimental evaluation of a
research codebase for financial time-series generation based on
Schrodinger-bridge ideas, jump-diffusion stochastic differential equations, and
modern neural baselines. The central objective is to generate realistic market
paths that preserve distributional shape, temporal dependence, cross-feature
correlation, and stylized facts such as heavy tails, volatility clustering, and
leverage effects. The implemented model family, called JD-SBTS, decomposes the
problem into jump detection, local volatility calibration, neural drift
estimation, and Euler-Maruyama simulation. A feedback variant introduces an
explicit stress process to amplify volatility after jumps. The benchmark also
includes Light Schrodinger Bridge variants, TimeGAN, an autoregressive RNN, and
a causal autoregressive Transformer. On the current QQQ stock benchmark
(60-day windows, six OHLCV-style features, 512 evaluation samples), the base
JD-SBTS model obtains the best aggregate rank, while the feedback variant
JD-SBTS-F gives the best stylized-fact preservation. The results show that different
evaluation criteria reward different modeling assumptions: autoregressive
neural models reproduce short-horizon conditional dynamics well, whereas
jump-diffusion bridge models better expose interpretable mechanisms for
discontinuities and volatility response.
\end{abstract}

\tableofcontents
\newpage

\section{Introduction}

Generating realistic financial time series is a difficult problem because market
data are not well described by a single Gaussian diffusion. Even at daily
frequency, price and return series display abrupt moves, non-Gaussian tails,
serial dependence in volatility, asymmetric responses to shocks, and
cross-sectional dependence between features such as open, high, low, close,
adjusted close, and volume. These properties are important in downstream tasks:
synthetic data used for risk analysis, stress testing, strategy backtesting, or
forecasting must not only look plausible path by path, but must also preserve
distributional and dynamic structure.

The project studies this problem through the lens of Schrodinger bridges. A
Schrodinger bridge can be viewed as a stochastic transport problem: among
stochastic processes close to a reference dynamics, find a process whose
endpoint marginals match observed distributions. This formulation is attractive
for time-series generation because it combines probabilistic dynamics with
distribution matching. However, financial time series also contain jumps and
volatility feedback, which are not naturally captured by a simple Brownian
reference. The codebase therefore develops jump-diffusion Schrodinger Bridge
Time Series models, abbreviated JD-SBTS, and compares them with neural and
potential-based alternatives.

The main contributions of the experiment are:
\begin{itemize}
    \item a unified implementation of several time-series generators behind the
    same \texttt{fit}/\texttt{generate} interface;
    \item a jump-diffusion SBTS model with static jump detection, local
    volatility calibration, LSTM drift estimation, and fast Euler generation;
    \item feedback and neural-jump extensions that model post-jump volatility
    amplification and data-dependent jump intensity;
    \item LightSB, TimeGAN, diffusion, RNN, and Transformer baselines in the
    same benchmark pipeline;
    \item a fair-comparison notebook protocol using shared training windows,
    shared evaluation subsets, and explicit generation-protocol labels.
\end{itemize}

The report focuses on the current maintained pipeline in \texttt{main.py},
\texttt{models/factory.py}, and \texttt{main\_pipeline.ipynb}. The numerical
tables use the latest saved output of the benchmark notebook for the QQQ stock
dataset. Older result figures for Merton and Google-style experiments remain in
the repository and are useful qualitative artifacts, but the main quantitative
discussion below is based on the current QQQ benchmark.

\section{Background}

\subsection{Financial Time-Series Generation}

Classical financial models often start from diffusion dynamics,
\begin{equation}
    dX_t = \mu(X_t,t)\,dt + \sigma(X_t,t)\,dW_t,
\end{equation}
where $\mu$ is a drift, $\sigma$ is a volatility coefficient, and $W_t$ is
Brownian motion. This model is analytically convenient, but it is too smooth for
many financial series. Empirical paths contain rare but large moves, and the
variance of future returns often increases after a large shock. Merton's
jump-diffusion model addresses part of this issue by adding a jump component:
\begin{equation}
    dX_t = \mu(X_t,t)\,dt + \sigma(X_t,t)\,dW_t + dJ_t,
\end{equation}
where $J_t$ is commonly driven by a Poisson process with random jump sizes.

For generative modeling, the goal is not only to specify a plausible stochastic
differential equation, but also to learn its components from data and generate
new paths whose distribution resembles the empirical one. Neural generative
models such as TimeGAN, diffusion models, recurrent networks, and Transformers
provide flexible alternatives. Their strength is representation learning, while
their weakness is often interpretability: a good sample may not reveal which
financial mechanism produced it. The JD-SBTS family tries to keep a mechanistic
SDE structure while still using neural components where they are useful.

\subsection{Schrodinger Bridges}

Let $R$ be a reference stochastic process on path space and let $\mu_0$ and
\(\mu_1\) be two target marginal distributions. The static form of the
Schrodinger bridge problem can be written as
\begin{equation}
    \pi^\star
    =
    \arg\min_{\pi \in \Pi(\mu_0,\mu_1)}
    \mathrm{KL}(\pi \Vert R_{01}),
\end{equation}
where $\Pi(\mu_0,\mu_1)$ is the set of couplings with the required marginals
and $R_{01}$ is the endpoint coupling induced by the reference process. The
dynamic form searches for a path measure close to $R$ while matching endpoint
constraints. For a Brownian reference, this can be interpreted as an
entropy-regularized optimal transport problem.

In practice, exact Schrodinger bridge solvers may be expensive in high
dimension. Light Schrodinger Bridge methods simplify the problem by learning
potential functions with a tractable parameterization. In this project, the
window-level LightSB baseline treats each full time-series window as one
high-dimensional point. A path-level LightSB variant also exists: it learns
one-step bridges on augmented states $[X_t,t]\rightarrow[X_{t+1},t+\Delta t]$.
This distinction matters because a flattened-window bridge is unconditional
with respect to the initial value, while a path-level bridge can be conditioned
on a real initial state.

\section{Implemented Models}

All active models are created through \texttt{models/factory.py}. Each model
implements the same high-level API:
\begin{equation}
    \texttt{model.fit(data, time\_grid)}, \qquad
    \texttt{model.generate(n\_samples, n\_steps, x0)}.
\end{equation}
This shared interface is essential for fair benchmarking because training,
generation, and evaluation can be orchestrated by the same loop.

\subsection{JD-SBTS: Jump-Diffusion SDE Framework}

The JD-SBTS (Jump-Diffusion Schrödinger Bridge Time Series) model is the
cornerstone of this research. It is formulated as a stochastic differential
equation (SDE) with jumps:
\begin{equation}
    dX_t = \mu(X_t, t)dt + \sigma(X_t, t)dW_t + dJ_t,
\end{equation}
where $W_t$ is a standard Wiener process and $J_t$ is a compound Poisson process.
The model follows a \textit{decoupled training} philosophy to ensure that each
physical component is identified without interference.

\paragraph{Jump Detection and Data Purification.}
To prevent volatility estimators from "double counting" variance during extreme
moves, the model first identifies jumps using a robust thresholding method:
$|\Delta X_t| > k \cdot \text{Med}(|\Delta X|) \cdot \sqrt{\Delta t}$. Once
flagged, jumps are removed, and the data is linearly interpolated to create a
\textit{purified} continuous-path dataset for volatility and drift calibration.

\paragraph{Local Volatility Calibration.}
The diffusion coefficient $\sigma(X_t, t)$ is estimated from purified data. We
implement a local volatility surface $\sigma(x,t)$ that captures the "volatility
smile" and state-dependent heteroscedasticity using kernel-based non-parametric
regression or state-grid interpolation.

\paragraph{Neural Drift Estimation.}
The drift term $\mu(X_t, t)$ represents the conditional expectation of the next
move. While traditional models use linear or kernel drift, JD-SBTS employs an
LSTM (Long Short-Term Memory) network to act as a deep drift estimator:
\begin{equation}
    \widehat{\mu}(X_t, t) = \text{LSTM}_{\theta}(\{X_s\}_{s \le t}).
\end{equation}
This allows the model to capture long-range dependencies and non-linear mean
reversion that are common in financial regimes.

\paragraph{Euler-Maruyama Simulation.}
For path generation, we use a discretized Euler scheme. A critical addition is
\textit{drift dampening}, where the neural drift is scaled by a factor $\eta \in (0,1)$
to ensure long-term stability and prevent the model from over-reacting to
historical noise during synthesis.

\subsection{Feedback JD-SBTS (JD-SBTS-F)}

The \texttt{jd\_sbts\_f} variant introduces an explicit mechanism for
\textit{volatility clustering}. We model a latent stress process $S_t$ that
evolves according to:
\begin{equation}
    dS_t = -\kappa S_t dt + \gamma |dJ_t|,
\end{equation}
where $\kappa$ is the decay rate and $\gamma$ is the sensitivity to jumps. The
instantaneous volatility is then modulated as $\sigma_{f}(X_t, t) = \sigma(X_t, t)
\sqrt{1 + S_t}$. This ensures that large jumps trigger a transient period of
high volatility, matching the empirical "leverage effect" and clustering facts.

\subsection{Light Schrödinger Bridge (LightSB)}

Based on the recent work of Korotin et al. (ICLR 2024), LightSB provides a
lightweight alternative to traditional Schrödinger Bridge solvers. Instead of
solving the SDE directly, it learns a pair of Schrödinger potentials $(v, \widehat{v})$
parameterized as a sum-exp of quadratic functions:
\begin{equation}
    v_{\theta}(x) = \sum_{k=1}^K \alpha_k \exp\left( -\frac{1}{2\epsilon} \|x - \mu_k\|_{S_k}^2 \right).
\end{equation}
The training objective is to maximize the log-likelihood of matching the target
marginal distribution $\mu_1$ from the source $\mu_0$ under the reference
measure. LightSB is particularly efficient because its conditional maps and
drifts can be computed analytically from the learned potentials. We implement
two versions: a window-level bridge (\texttt{lightsb}) and a sequential
path-level bridge (\texttt{lightsb\_path}).

\subsection{Deep Learning Baselines}

\paragraph{TimeGAN.}
A Generative Adversarial Network architecture specifically tailored for time
series. It consists of four integrated networks: an embedding network, a
recovery network, a generator, and a discriminator. Unlike standard GANs, it
includes a \textit{supervised loss} that encourages the generator to capture the
step-wise conditional distributions, preventing mode collapse.

\paragraph{Autoregressive RNN and Transformer.}
The \texttt{rnn} baseline uses a standard GRU/LSTM architecture for next-step
prediction. The \texttt{transformer\_ar} utilizes the multi-head attention
mechanism with causal masking to capture global dependencies across the entire
window. Both are trained using Maximum Likelihood Estimation (MLE) and generate
paths by recursive sampling.

\subsection{Model Summary and Generation Protocol}

\begin{table}[H]
\centering
\caption{Main models and generation protocols used by the benchmark.}
\begin{tabular}{lll}
\toprule
Model & Main mechanism & Protocol in benchmark \\
\midrule
\texttt{jd\_sbts} & Jump-diffusion SBTS & Shared real initial state \\
\texttt{jd\_sbts\_f} & JD-SBTS with stress feedback & Shared real initial state \\
\texttt{lightsb} & Flattened-window potential bridge & Shared seed only \\
\texttt{rnn} & Autoregressive recurrent network & Shared real prefix, length 20 \\
\texttt{timegan} & Adversarial recurrent generator & Shared seed only \\
\texttt{transformer\_ar} & Causal autoregressive Transformer & Shared real prefix, length 20 \\
\bottomrule
\end{tabular}
\end{table}

\section{Experimental Setup}

\subsection{Datasets}

The notebook supports both simulated and real datasets. The simulated datasets
are Merton jump-diffusion, standard Ornstein-Uhlenbeck, and high-frequency
Ornstein-Uhlenbeck paths. The main quantitative benchmark discussed here uses
real stock data downloaded through Yahoo Finance.

For the current run, the selected dataset is QQQ:
\begin{itemize}
    \item ticker: QQQ;
    \item requested date range: 2004-01-01 to 2019-12-31;
    \item actual date range in the saved data: 2004-01-02 to 2019-12-31;
    \item valid rows: 4027;
    \item features: High, Low, Open, Close, Adj Close, Volume;
    \item window length: 60;
    \item number of processed windows: 3968;
    \item normalization: base-one normalization within each window.
\end{itemize}

The split is 70\% train, 15\% validation, and 15\% test, with shuffled split
indices stored for reproducibility. The resulting shapes are:
\begin{equation}
    \text{train}: (2777,60,6), \quad
    \text{validation}: (595,60,6), \quad
    \text{test}: (596,60,6).
\end{equation}
The evaluation subset uses the first 512 windows from the test pool.

\subsection{Training Configuration}

The notebook fixes the random seed at 42 and applies model-specific training
overrides. The main values are:
\begin{itemize}
    \item LSTM drift epochs: 50;
    \item TimeGAN epochs: 80;
    \item LightSB epochs: 50;
    \item RNN epochs: 75, hidden dimension 128, context length 20;
    \item Transformer epochs: 75, model dimension 128, three layers, four
    attention heads, context length 20;
    \item generated/evaluated samples: 512.
\end{itemize}

The comparison is not identical conditioning for every model, because the model
interfaces differ. Instead, the notebook makes the conditioning explicit.
JD-SBTS models receive the same real initial states. Autoregressive models
receive the same real prefix windows of length 20 and complete the remaining
trajectory. Window-level and adversarial models are unconditional with respect
to the specific evaluation initial state, so they share the reset random seed
and the same evaluation subset.

\subsection{Metrics}

The evaluation uses complementary metrics rather than a single score. Lower
values are better for all metric values in the main ranking table.

\paragraph{Wasserstein distance.}
The average one-dimensional Wasserstein distance compares empirical return
distributions feature by feature and averages over all features. It measures
whether generated increments have the same marginal distribution as real
increments without mixing adjacent feature columns into artificial differences.

\paragraph{Autocorrelation MSE.}
The ACF error computes return autocorrelation functions for every
sample-feature series, averages them, and compares the real and generated
average ACF curves up to a fixed maximum lag. It penalizes incorrect temporal
dependence across all features.

\paragraph{Correlation distance.}
The correlation metric computes the Frobenius distance between real and
generated feature-correlation matrices. It is important for multi-feature stock
windows because the six input channels are strongly dependent.

\paragraph{Discriminative scores.}
The pipeline records two post-hoc classifier diagnostics. The GRU
discriminator tests sequential distinguishability and also records an
effective score that accounts for prediction-balance bias. The 1D-CNN
discriminator tests whether local temporal patterns make generated windows
separable from real windows. Both raw scores are defined as
\(|\mathrm{accuracy}-0.5|\), so lower values mean that the classifier is closer
to random guessing. These diagnostics should still be read carefully: they
measure realism from specific classifier perspectives, not calibrated
statistical distances.

\paragraph{Predictive score.}
A predictor is trained on generated windows and evaluated on real windows. This
tests whether generated sequences contain transferable forecasting structure.
When no target column is specified, the current implementation predicts all
features at the final step and reports the MAE.

\paragraph{Stylized facts.}
The notebook computes relative errors for volatility clustering, fat tails, and
leverage effect. These quantities are computed feature by feature from returns
and then averaged, so the reported stylized score now reflects the full
multi-feature stock window rather than only a single selected column.

\section{Results}

\subsection{Runtime and Primary Metrics}

Table~\ref{tab:main-results} reports the current saved QQQ benchmark values for
the seven models that appear in the latest metric table. The base JD-SBTS model
has the best aggregate rank because it combines the best ACF error, a strong
predictive score, and good stylized-fact preservation. The RNN is a close
second and remains strong on short-horizon prefix completion. The Transformer
has the best Wasserstein distance and the best correlation distance, but its
stylized-fact score is weaker. JD-SBTS-F has the best predictive score and the
best stylized-fact score, although its marginal Wasserstein and correlation
metrics are weaker in this run.

\begin{table}[H]
\centering
\caption{Main QQQ benchmark results. Lower is better for Wasserstein distance,
ACF MSE, predictive score, correlation distance, stylized score, and overall
rank score.}
\label{tab:main-results}
\resizebox{\textwidth}{!}{%
\begin{tabular}{lrrrrrrr}
\toprule
Model & Train time (s) & Gen. time (s) & Wasserstein & ACF MSE
& Predictive & Corr. dist. & Overall rank \\
\midrule
\texttt{jd\_sbts} & 156.52 & 0.45 & 0.103038 & 0.000992 & 0.274136 & 3.913000 & 2.967 \\
\texttt{rnn} & 8.09 & 0.13 & 0.031686 & 0.001911 & 0.274640 & 0.468194 & 3.033 \\
\texttt{transformer\_ar} & 16.16 & 0.26 & 0.014647 & 0.006709 & 0.283070 & 0.179678 & 3.233 \\
\texttt{jd\_sbts\_f} & 152.54 & 0.29 & 0.172726 & 0.005617 & 0.269187 & 4.796639 & 4.200 \\
\texttt{lightsb} & 2.92 & 0.10 & 0.146802 & 0.002044 & 0.318087 & 3.347431 & 4.300 \\
\texttt{lightsb\_path} & 146.60 & 0.10 & 0.039541 & 0.008194 & 0.300270 & 3.549825 & 4.367 \\
\texttt{timegan} & 103.80 & 0.02 & 0.066906 & 0.159717 & 1.665162 & 3.672485 & 5.900 \\
\bottomrule
\end{tabular}}
\end{table}

Several observations are immediate. First, pure runtime is not aligned with
metric quality. LightSB is very fast to train in this configuration, but its
window-level formulation is still weaker on the overall ranking. Second, the
autoregressive baselines are strong on this 60-day stock-window task: the RNN
is close to JD-SBTS overall, while the Transformer gives the closest terminal
marginal distribution and feature-correlation structure. Third, JD-SBTS
variants are slower to train because they combine calibration, drift learning,
and simulation components, but this run shows that their mechanistic structure
can be competitive with purely neural sequence models.

\subsection{Discriminative Diagnostics}

Table~\ref{tab:disc} shows the saved GRU discriminative diagnostics from the
current benchmark output. The updated pipeline additionally records a 1D-CNN
discriminative score under \texttt{cnn\_discriminative\_score}; the full table
should be regenerated after rerunning the notebook. The diagnostic uses 512
real and 512 generated windows, repeated 10 times with an 80/20 balanced split.
Each repeat trains on 409 real and 409 generated windows and tests on 103 real
and 103 generated windows.

\begin{table}[H]
\centering
\caption{Saved GRU discriminative diagnostics. The score is derived from a
post-hoc GRU classifier and is defined as $|\mathrm{accuracy}-0.5|$; lower
values indicate that the classifier has more difficulty separating real and
generated windows.}
\label{tab:disc}
\begin{tabular}{lrrrr}
\toprule
Model & Score & Std. & Accuracy & Predicted-real rate \\
\midrule
\texttt{lightsb} & 0.296117 & 0.019659 & 0.796117 & 0.499029 \\
\texttt{transformer\_ar} & 0.332039 & 0.091417 & 0.832039 & 0.499029 \\
\texttt{jd\_sbts} & 0.371845 & 0.125656 & 0.871845 & 0.580583 \\
\texttt{timegan} & 0.450485 & 0.143722 & 0.950485 & 0.450485 \\
\texttt{lightsb\_path} & 0.458252 & 0.017664 & 0.958252 & 0.494175 \\
\texttt{rnn} & 0.482524 & 0.016704 & 0.982524 & 0.504854 \\
\texttt{jd\_sbts\_f} & 0.488350 & 0.008464 & 0.988350 & 0.502913 \\
\bottomrule
\end{tabular}
\end{table}

The discriminative table should not be interpreted alone. For instance,
LightSB obtains the lowest classifier-based score, but it does not dominate the
primary statistical metrics. TimeGAN has a moderate discriminative score while
its predictive score is poor and its ACF error is large. This indicates that a
discriminator can be insensitive to some properties that matter for forecasting
transfer or temporal dependence. Conversely, JD-SBTS-F has the best
stylized-fact score but is easy for the classifier to separate from real data.

\subsection{Stylized-Fact Preservation}

Table~\ref{tab:stylized} reports relative errors for the three stylized facts.
JD-SBTS-F has the best average stylized score, followed by the base JD-SBTS.
This is consistent with the model design: explicit jumps, local volatility, and
feedback stress are intended to capture tail behavior and volatility response,
whereas autoregressive neural models mainly optimize conditional next-step
dynamics.

\begin{table}[H]
\centering
\caption{Stylized-fact relative errors on generated returns. Lower is better.}
\label{tab:stylized}
\begin{tabular}{lrrrr}
\toprule
Model & Vol. clustering & Fat tails & Leverage effect & Stylized score \\
\midrule
\texttt{jd\_sbts\_f} & 0.028114 & 0.345227 & 0.095700 & 0.156347 \\
\texttt{jd\_sbts} & 0.360012 & 0.533939 & 0.346011 & 0.413321 \\
\texttt{lightsb\_path} & 0.511614 & 0.438253 & 0.964585 & 0.638150 \\
\texttt{lightsb} & 0.766930 & 0.983955 & 1.003505 & 0.918130 \\
\texttt{timegan} & 0.905866 & 0.981688 & 0.959081 & 0.948878 \\
\texttt{transformer\_ar} & 0.689446 & 3.352891 & 0.240817 & 1.427718 \\
\texttt{rnn} & 0.910866 & 4.042890 & 0.085005 & 1.679587 \\
\bottomrule
\end{tabular}
\end{table}

The feedback model now clearly improves stylized-fact preservation. Its
volatility-clustering and leverage-effect errors are especially small relative
to the other models. At the same time, it remains weaker on Wasserstein and
correlation distance, which shows that matching financial stylized facts and
matching full multivariate window distributions are related but distinct
objectives.

\subsection{Qualitative Path Comparisons}

Figure~\ref{fig:paths} compares real and generated close-price trajectories for
three representative models. The left panel of each subfigure shows matched
individual windows, and the right panel shows the 10--90\% distribution band and
mean trajectory.

\begin{figure}[H]
\centering
\begin{subfigure}{0.95\textwidth}
    \centering
    \includegraphics[width=\textwidth]{\detokenize{notebook_outputs/benchmark_main_pipeline/visual_comparisons/real_vs_generated_rnn_Close.png}}
    \caption{\texttt{rnn}: strong aggregate rank and smooth prefix-conditioned completions.}
\end{subfigure}
\vspace{0.4em}
\begin{subfigure}{0.95\textwidth}
    \centering
    \includegraphics[width=\textwidth]{\detokenize{notebook_outputs/benchmark_main_pipeline/visual_comparisons/real_vs_generated_jd_sbts_f_Close.png}}
    \caption{\texttt{jd\_sbts\_f}: jump-diffusion dynamics with feedback stress.}
\end{subfigure}
\vspace{0.4em}
\begin{subfigure}{0.95\textwidth}
    \centering
    \includegraphics[width=\textwidth]{\detokenize{notebook_outputs/benchmark_main_pipeline/visual_comparisons/real_vs_generated_transformer_ar_Close.png}}
    \caption{\texttt{transformer\_ar}: best terminal Wasserstein distance in the current table.}
\end{subfigure}
\caption{Real versus generated close-feature paths from the saved visual
comparison artifacts.}
\label{fig:paths}
\end{figure}

The visual comparisons support the quantitative results. The autoregressive
models follow the upward drift and conditional continuation implied by the
prefix. The JD-SBTS-F samples are less tightly tied to the prefix trend but
preserve a stochastic path-generation interpretation with volatility response.
This is a central trade-off of the project: neural autoregressive baselines can
be very competitive on finite-window benchmarks, while bridge/SDE models
provide a more explicit stochastic structure for scenario generation.

\section{Discussion}

\subsection{What Worked}

The unified factory design is one of the strongest engineering outcomes. It
allows very different models to be evaluated with the same pipeline and reduces
the risk that result differences come from inconsistent data handling. The
notebook also explicitly records generation protocols. This is important
because conditioning differs by model class: comparing a prefix-conditioned RNN
with an unconditional window-level bridge without stating the protocol would be
misleading.

The benchmark confirms that simple autoregressive models are hard to beat on
short stock windows. The RNN is not the most elaborate model, but it receives a
real 20-step prefix and learns the next-step dynamics well. This gives it low
ACF error, a strong predictive score, and the second-best aggregate rank. The
Transformer improves terminal Wasserstein and feature-correlation distance,
suggesting that attention helps match marginal and cross-feature structure, but
it does not improve all temporal or stylized metrics.

The JD-SBTS models are valuable because they target mechanisms rather than only
black-box sequence likelihood. The base model has the best aggregate rank in
the current table, while the feedback variant has the best stylized-fact score.
This is important for financial applications where tail risk, volatility
clustering, and leverage-like behavior may matter more than a single aggregate
rank.

\subsection{Limitations}

Several limitations remain.

First, the current quantitative table is a single benchmark run on QQQ windows.
The repository supports Merton, OU, high-frequency OU, and other stock tickers,
but a robust empirical conclusion requires repeated runs across datasets and
seeds. A single random split can favor models whose inductive bias matches the
selected window distribution.

Second, the generation protocols are fair but not identical. Prefix-conditioned
models have more information about each evaluation trajectory than
unconditional models. This is acceptable when stated explicitly, but future
tables should separate unconditional generation, initial-state conditioning,
and prefix completion as different tasks.

Third, the current LightSB window baseline is disadvantaged by flattening the
entire path. A 60-by-6 window becomes a 360-dimensional vector, and temporal
structure is only implicit. The path-level LightSB variant is a better
conceptual match for time series and should be included systematically in the
final quantitative tables.

Fourth, discriminative scores are sensitive to classifier capacity and training
details. Reporting both GRU and 1D-CNN variants reduces dependence on a single
classifier architecture, but low classifier-based scores still do not imply
good predictive or stylized-fact performance. They should be treated as
diagnostics among several, not as the main target.

\subsection{Next Improvements}

The most direct next step is to run the same full benchmark for all 11
registered models and save the metric tables as CSV files. This would make the
report easier to reproduce and avoid relying on notebook display outputs.
Additional improvements include:
\begin{itemize}
    \item repeated seeds with mean and confidence intervals;
    \item a separated task taxonomy: unconditional generation, $x_0$-conditioned
    generation, and prefix completion;
    \item better calibration of feedback parameters $\kappa$ and $\gamma$;
    \item explicit parameter counts, training time, and generation time for all
    models;
    \item additional market assets and crisis periods to test tail robustness;
    \item direct comparison between \texttt{lightsb} and
    \texttt{lightsb\_path} on the same metrics.
\end{itemize}

\section{Conclusion}

The project implements a broad and extensible framework for financial
time-series generation. The JD-SBTS family introduces jump-diffusion
Schrodinger-bridge-inspired generation with interpretable components: jump
detection, local volatility calibration, neural drift estimation, and
feedback-driven volatility amplification. The benchmark demonstrates that no
single model dominates every criterion. On the current QQQ benchmark, the base
JD-SBTS model has the best aggregate rank, the Transformer has the best terminal
Wasserstein distance and correlation distance, and JD-SBTS-F has the best
stylized-fact score. This split of winners is important: it shows that the
choice of model should depend on whether the priority is distribution matching,
prefix completion, interpretable dynamics, or financial stylized facts.

The main lesson is that model quality depends on the intended use case. If the
task is short-horizon prefix completion, autoregressive neural baselines are
strong and efficient. If the task is interpretable scenario generation with
jumps, tails, and volatility mechanisms, the JD-SBTS approach is more aligned
with the financial structure of the data. Future work should therefore evaluate
models by task type and extend the benchmark across more datasets, seeds, and
market regimes.

\begin{thebibliography}{9}

\bibitem{schrodinger1931}
E. Schrodinger.
\newblock Uber die Umkehrung der Naturgesetze.
\newblock \emph{Sitzungsberichte der Preussischen Akademie der Wissenschaften},
1931.

\bibitem{merton1976}
R. C. Merton.
\newblock Option pricing when underlying stock returns are discontinuous.
\newblock \emph{Journal of Financial Economics}, 1976.

\bibitem{leonard2014}
C. Leonard.
\newblock A survey of the Schrodinger problem and some of its connections with
optimal transport.
\newblock \emph{Discrete and Continuous Dynamical Systems A}, 2014.

\bibitem{timegan2019}
J. Yoon, D. Jarrett, and M. van der Schaar.
\newblock Time-series generative adversarial networks.
\newblock \emph{NeurIPS}, 2019.

\bibitem{transformer2017}
A. Vaswani et al.
\newblock Attention is all you need.
\newblock \emph{NeurIPS}, 2017.

\bibitem{ddpm2020}
J. Ho, A. Jain, and P. Abbeel.
\newblock Denoising diffusion probabilistic models.
\newblock \emph{NeurIPS}, 2020.

\bibitem{lightsb2024}
A. Korotin et al.
\newblock Light Schrodinger Bridge.
\newblock \emph{ICLR}, 2024.

\end{thebibliography}

\end{document}
